\documentclass[11pt]{article}

\usepackage[utf8]{inputenc}
\usepackage[T1]{fontenc}
\usepackage{amsmath}
\usepackage{amssymb}
\usepackage{booktabs}
\usepackage{graphicx}
\usepackage{hyperref}
\usepackage{microtype}
\usepackage{enumitem}

\IfFileExists{neurips_2025.sty}{
  \usepackage{neurips_2025}
}{
  \usepackage{geometry}
  \geometry{margin=1in}
}

\title{Replacing Transformer FFNs with Localized Recurrent Memory}
\author{
Liu Xiao \\
\texttt{liu.xiao.in@gmail.com}
\And
GPT-5.4 xhigh
}
\date{}

\begin{document}
\maketitle

\begin{abstract}
Transformer feed-forward networks are usually treated as token-local nonlinear
blocks, but they can also be viewed as a form of static memory stored in
weights. We study a narrow replacement question: can the FFN slot in a
decoder-only Transformer be replaced by a localized recurrent memory module
while keeping self-attention fixed? This isolates dynamic memory from full
backbone recurrence. We first screen a family of recurrent FFN replacements,
including token-conditioned readout, matrix-state memory, multiscale decay,
predictive-corrective updates, stable dynamics, sparse writes, and auxiliary
state prediction. On synthetic memory tasks, the plain Recurrent-FFN already
improves over Mamba2 and softmax attention baselines in our small pilot,
reaching 15.8\% associative-recall accuracy versus 10.2\% and 7.0\%,
respectively. On byte-level WikiText-2, several recurrent variants beat a
full-width SwiGLU Transformer at 120 steps, and stable/readout variants are the
best fair small-model performers at 200 steps. We then make a deliberately
accuracy-first move: an \textbf{Ultra} compound recurrent FFN that mixes a
readout memory, a stable recurrent memory, and a widened local SwiGLU branch.
In a two-seed, 400-step matched-capacity comparison on byte-level WikiText-2,
Ultra-384 (1,294,976 parameters) reaches 3.2544 and 3.2544 bits-per-byte,
versus 3.4744 and 3.4384 for a Transformer-SwiGLU with 1,294,880 parameters.
The systems story remains poor in the current Python reference
implementation: the matched-capacity Transformer reaches roughly
204k--281k tokens/sec at sequence lengths 64--256 on Apple MPS, while
Ultra-384 reaches only about 6.6k--6.7k. We therefore interpret localized
recurrent FFNs as a credible small-scale quality and memory win, with the
strongest current result coming from an accuracy-first compound design rather
than the original strict parameter-matched replacement.
\end{abstract}

\section{Introduction}
Modern recurrent and state-space models have revived interest in replacing
quadratic attention with learned state updates. Earlier attempts to inject
recurrence into Transformer-style computation include Universal Transformers
and direct Transformer-to-RNN conversion work \citep{dehghani2019universal,
kasai2021t2r}. More recent matrix-state recurrence, Linear Recurrent Units,
RWKV-style recurrence, selective state spaces, and hybrid recurrent systems all
suggest that carefully designed state can recover more sequence structure than
older recurrent architectures were able to preserve
\citep{orvieto2023lru, rwkv2023, eaglefinch2024, mamba2023,
mamba2ssd2024, xlstm2024, m2rnn2026}. Recent models such as RetNet,
DeltaNet, Griffin, and Gated Delta Networks further sharpen this design space
by combining recurrent or state-space updates with retention, local attention,
delta-rule memory updates, or sequential linear-time mixing
\citep{retnet2023, yang2024deltanet, griffin2024, gateddeltanet2024}. Hybrid
designs further suggest that strong sequence models may emerge not from a
single replacement of attention, but from careful division of labor between
attention, recurrence, and explicit memory \citep{jamba2024, samba2024,
hymba2024}. At the same time, full backbone
replacement makes it difficult to disentangle which parts of Transformer
quality come from self-attention and which come from the feed-forward layers.

This paper studies a narrower question:
\begin{quote}
Can we replace only the Transformer FFN with localized recurrent memory while
keeping attention fixed?
\end{quote}
That framing is useful for two reasons. First, it isolates one component of the
Transformer block instead of turning the experiment into a complete backbone
change. Second, it gives a clean bridge from Transformer FFNs, which can be
interpreted as a kind of static memory stored in parameters, to modern
token-conditioned recurrent memory. That bridge is not only intuitive but also
supported by prior mechanistic work: Transformer feed-forward layers have been
interpreted as key-value memories stored in weights \citep{geva2020ffnmem},
knowledge-bearing FFN units have been localized in pretrained models
\citep{dai2021knowledge}, direct factual editing results often operate on MLP
weights \citep{meng2022rome}, and recent FFN-centric memory work has explored
making such memory more explicit and controllable \citep{memoryllmffn2026}.
Our paper asks whether a localized
recurrent substitute can play a similar role while adding token-conditioned
dynamics.

Our starting point is simple. We keep the standard decoder-only residual
layout,
\[
x \rightarrow \text{attention} \rightarrow \text{residual} \rightarrow
\text{FFN replacement} \rightarrow \text{residual},
\]
and change only the FFN slot. This creates a localized recurrent memory block
inside an otherwise standard Transformer. We then ask five progressively more
refined questions:
\begin{enumerate}[leftmargin=1.25em]
  \item Does a plain recurrent FFN help memory-sensitive behavior at all?
  \item If language-model quality lags, is the missing ingredient token-
  conditioned memory readout?
  \item Does the recurrent block need to recover some of the token-local
  nonlinear mixing of a normal Transformer MLP?
  \item Which recurrent inductive bias best uses the FFN slot: matrix memory,
  multiscale decay, predictive correction, stable dynamics, sparse writes, or
  auxiliary prediction?
  \item If the fair drop-in variants are still not strong enough, can an
  accuracy-first compound recurrent FFN beat a matched-capacity Transformer?
\end{enumerate}

The resulting story is narrower than a general ``RNNs beat Transformers''
claim, but more scientifically legible. Our experiments show that the plain
localized recurrent block is already a memory-sensitive win on synthetic tasks,
yet initially slightly worse on byte-level WikiText-2. A wider method screen
then shows that several recurrent FFNs beat a small strong SwiGLU baseline. The
most important update in this paper is that the story no longer ends there:
once we introduce an accuracy-first compound recurrent FFN and compare it to a
matched-capacity Transformer, the recurrent FFN can achieve a much larger BPB
gain in this small benchmark.

\paragraph{Contributions.}
\begin{itemize}[leftmargin=1.25em]
  \item We define a localized recurrent replacement for the Transformer FFN
  that keeps self-attention, residual structure, and normalization fixed.
  \item We screen a broad family of localized recurrent FFN variants: readout,
  hybrid, matrix-state, multiscale, predictive-corrective, stable,
  sparse-write, and auxiliary-prediction designs.
  \item We show that the plain localized recurrent FFN improves small
  synthetic memory tasks even before it improves language-model quality.
  \item We show that several recurrent variants beat a full-width SwiGLU
  Transformer at 120 steps, and that stable/readout variants are the strongest
  small matched-capacity performers at 200 steps.
  \item We introduce an accuracy-first \emph{Ultra} compound recurrent FFN and
  show that, in a two-seed 400-step matched-capacity comparison, it beats a
  larger Transformer-SwiGLU by roughly 0.20 BPB on byte-level WikiText-2.
  \item We report the caveat clearly: these localized recurrent FFNs remain far
  slower than the Transformer in the current Python reference implementation.
\end{itemize}

\section{Method}
\subsection{Localized Recurrent FFN Replacement}
We start from a standard decoder-only Transformer block with self-attention and
residual connections \citep{vaswani2017attention}. The attention path is left
unchanged. Only the post-attention FFN is replaced.

Given normalized token features $u_t \in \mathbb{R}^{d}$, the plain
Recurrent-FFN updates a per-layer recurrent state $s_t \in \mathbb{R}^{m}$:
\begin{align}
f_t &= \tanh(W_f u_t), \\
i_t &= \sigma(W_i u_t), \\
v_t &= \mathrm{silu}(W_v u_t), \\
s_t &= f_t \odot s_{t-1} + i_t \odot v_t, \\
y_t &= W_o \, \mathrm{silu}(s_t).
\end{align}
The output $y_t$ is projected back to model width and added through the normal
residual path.

\subsection{Readout, Local Mixing, and Screened Recurrent Variants}
The first baseline removes most of the standard FFN's token-local expansion and
asks the recurrent state to do almost all of the work. To recover some of that
capacity, we first add token-conditioned readout and then screen several
physics- and memory-inspired recurrent update rules.

\paragraph{Recurrent-FFN-Readout.}
This variant adds a token-conditioned readout vector
\begin{equation}
q_t = W_q u_t,
\end{equation}
and exposes recurrent memory through an elementwise interaction:
\begin{equation}
y_t = W_o \, \mathrm{silu}(q_t \odot s_t).
\end{equation}
This allows the current token to select which recurrent channels are relevant.

\paragraph{Recurrent-FFN-Hybrid.}
The simplest quality-oriented variant keeps the recurrent readout path but also
adds a local gated branch in the spirit of GLU/SwiGLU Transformer FFNs
\citep{shazeer2020glu}:
\begin{align}
r_t &= W_r \, \mathrm{silu}(q_t \odot s_t), \\
\ell_t &= W_\ell \left(\mathrm{silu}(U u_t) \odot \sigma(G u_t)\right), \\
y_t &= r_t + \ell_t.
\end{align}
This restores part of the token-local nonlinear mixing that a standard
Transformer FFN provides, while preserving a dynamic recurrent memory channel.

\paragraph{Screened recurrent update rules.}
On top of this readout and local-mixing template, we test six additional
recurrent inductive biases. The \textbf{matrix} variant writes a low-rank
matrix-valued state. The \textbf{multiscale} variant maintains several decay
timescales and learns a token-conditioned mixture over them. The
\textbf{predictive} variant uses a prediction-correction update reminiscent of
observer-style filtering. The \textbf{stable} variant applies damped,
token-conditioned rotations to paired recurrent channels. The \textbf{sparse}
variant enforces top-$k$ writes. The \textbf{auxiliary} variant adds a next-
feature prediction loss on the recurrent state.

\paragraph{Stable-dynamics variant.}
The strongest fair small-model result comes from the stable variant. It splits
the recurrent state into paired channels, applies token-conditioned rotation
and damping, and then reads the state through the same token-conditioned query
and local SwiGLU-like branch. Let the previous state be partitioned into
$(a_{t-1}, b_{t-1})$ pairs. For each pair, we apply
\begin{equation}
\mathrm{rot}(a, b; \omega) =
\left(a \cos \omega - b \sin \omega,\; a \sin \omega + b \cos \omega\right).
\end{equation}
The stable recurrent update is then
\begin{align}
\omega_t &= \pi \tanh(W_\omega u_t), \\
\tilde{s}_t &= \mathrm{rot}(s_{t-1}; \omega_t), \\
d_t &= \exp(-\mathrm{softplus}(W_d u_t)), \\
g_t &= \sigma(W_g u_t), \\
c_t &= \mathrm{silu}(W_c u_t), \\
s_t &= d_t \odot \tilde{s}_t + g_t \odot c_t.
\end{align}
The damped rotation is intended to preserve useful phase-like structure without
allowing the recurrent state to grow unstably.

\paragraph{Ultra compound variant.}
The strongest final result comes from a deliberately accuracy-first
\textbf{Ultra} variant. Ultra keeps the FFN slot localized, but allocates more
capacity and combines three subpaths:
\begin{enumerate}[leftmargin=1.25em]
  \item a readout memory path of the Recurrent-FFN-Readout form, producing
  $r_t$,
  \item a stable recurrent path of the form above, producing $h_t$,
  \item a widened local SwiGLU-style branch, producing $\ell_t$.
\end{enumerate}
The three paths are mixed with token-dependent coefficients:
\begin{align}
\alpha_t &= \mathrm{softmax}(W_{\text{mix}} u_t), \\
y_t &= \alpha_{t,1} r_t + \alpha_{t,2} h_t + \alpha_{t,3} \ell_t.
\end{align}
Unlike the earlier screen, Ultra is not intended to be a strict parameter-
matched drop-in replacement during the search phase. It is an accuracy-first
compound FFN. We only restore fairness at the final comparison step by matching
its total parameter count against a larger Transformer-SwiGLU control.

\section{Experimental Setup}
\paragraph{Language modeling.}
We train byte-level autoregressive models on WikiText-2 raw text
\citep{merity2017pointer}. All models are decoder-only and use sequence length
$128$, batch size $16$, and Apple MPS.

We report four regimes:
\begin{enumerate}[leftmargin=1.25em]
  \item a historical 80-step pilot with the original GELU Transformer baseline,
  \item a 120-step recurrent method screen at $d=64$, 3 layers, 4 heads, zero
  dropout, and learning rate $4 \times 10^{-4}$, with recurrent state sizes
  chosen to near-match the 223,040-parameter full-width SwiGLU baseline,
  \item a 200-step two-seed follow-up for the strongest fair recurrent
  candidates under the same small-model setup,
  \item an accuracy-first Ultra search with wider recurrent budgets and longer
  schedules, followed by matched-capacity controls against larger
  Transformer-SwiGLU models.
\end{enumerate}

\paragraph{Accuracy-first matched-capacity setting.}
The main final comparison uses Ultra-384 with $d=64$, 3 layers, 4 heads, and
state size 384, for 1,294,976 parameters. It is trained for 400 steps with
learning rate $3 \times 10^{-4}$, weight decay 0.02, zero dropout, and 40
warmup steps. The matched Transformer control uses $d=160$, 3 layers, 4 heads,
the standard full-width SwiGLU FFN, and 1,294,880 parameters. Among the
Transformer settings we tested, the best control uses 400 steps, learning rate
$4 \times 10^{-4}$, weight decay 0.05, and zero warmup.

\paragraph{Synthetic memory tasks.}
We also evaluate on small synthetic memory-sensitive tasks implemented in the
local benchmark harness: associative recall and signal-vs-noise interference.
The initial pilot uses sequence length $64$, hidden size $64$, 2 layers,
4 heads, batch size $8$, and 120 training steps.

\paragraph{Throughput.}
We measure forward throughput on Apple MPS at sequence lengths 64, 128, and
256. The throughput table compares the matched-capacity Transformer control
with Ultra-384. These are reference-implementation numbers rather than
optimized systems measurements, because the recurrent state updates are written
as explicit Python loops.

\section{Results}
\subsection{Initial Synthetic Memory Pilot}
The plain Recurrent-FFN baseline is already a positive result on memory-
sensitive tasks. Table~\ref{tab:memory} shows that it outperforms both Mamba2
and softmax attention in associative recall and signal-vs-noise accuracy in
our first small pilot.

\begin{table}[t]
  \centering
  \caption{Initial synthetic memory pilot at sequence length 64.}
  \label{tab:memory}
  \begin{tabular}{lccc}
    \toprule
    Task & Recurrent-FFN & Mamba2 & Softmax attention \\
    \midrule
    Associative recall accuracy & 15.82\% & 10.16\% & 7.03\% \\
    Signal-noise accuracy & 12.70\% & 8.59\% & 8.59\% \\
    \bottomrule
  \end{tabular}
\end{table}

This matters because attention is unchanged. The gain comes from replacing only
the FFN slot with a dynamic memory mechanism.

\subsection{Historical Small-Model Development}
The first matched-parameter WikiText-2 pilot used the original GELU Transformer
baseline and showed that the plain recurrent FFN was still slightly worse at 80
training steps: 5.178 BPB versus 5.144. We then changed the main baseline to a
full-width SwiGLU Transformer and reran a wider recurrent screen. Table
~\ref{tab:pilot120} shows the 120-step method screen.

\begin{table}[t]
  \centering
  \small
  \caption{120-step byte-level WikiText-2 method screen with zero dropout and
  learning rate $4 \times 10^{-4}$. Lower is better. The main baseline is the
  full-width Transformer-SwiGLU row; GELU and matched-SwiGLU rows provide
  context.}
  \label{tab:pilot120}
  \begin{tabular}{lcccc}
    \toprule
    Model & Params & Val BPB & Val ppl & Train sec \\
    \midrule
    Transformer-SwiGLU & 223,040 & 4.404 & 21.17 & 1.38 \\
    Transformer-GELU & 174,848 & 4.309 & 19.82 & 1.21 \\
    Transformer-SwiGLU-matched & 173,504 & 4.411 & 21.28 & 1.25 \\
    Recurrent-FFN (192) & 223,040 & 4.610 & 24.42 & 16.16 \\
    Readout (152) & 221,960 & 4.375 & 20.74 & 19.90 \\
    Hybrid (192) & 223,328 & 4.431 & 21.57 & 35.58 \\
    Matrix (384) & 212,987 & 4.386 & 20.91 & 40.93 \\
    Multiscale (216) & 222,152 & 4.368 & 20.65 & 33.24 \\
    Predictive (192) & 223,328 & 4.462 & 22.04 & 29.04 \\
    Stable (184) & 226,004 & 4.370 & 20.68 & 37.73 \\
    Sparse (192) & 223,328 & 4.428 & 21.52 & 33.80 \\
    Auxiliary (168) & 220,988 & 4.438 & 21.68 & 35.94 \\
    \bottomrule
  \end{tabular}
\end{table}

Four conclusions follow. First, optimization still matters a great deal for
localized recurrent FFNs: the plain recurrent baseline remains well behind.
Second, token-conditioned readout helps immediately. Third, once the search
space is widened, the old hybrid winner is no longer the best recurrent model.
Fourth, several recurrent variants beat the full-width Transformer-SwiGLU
baseline at 120 steps: Multiscale, Stable, Readout, and Matrix. However, the
smaller Transformer-GELU baseline remains the best short-budget Transformer in
this tiny setting, so the 120-step result is a SwiGLU-specific win rather than
yet a general FFN win.

Table~\ref{tab:pilot200} shows the 200-step two-seed follow-up for the best
fair recurrent candidates. The earlier hybrid winner disappears entirely.
Recurrent-FFN-Readout gives the best single-seed run in seed 7, while
Recurrent-FFN-Stable gives the best two-seed mean among the fair small models.

\begin{table}[t]
  \centering
  \small
  \caption{Two-seed 200-step byte-level WikiText-2 follow-up for the strongest
  fair recurrent candidates from the 120-step screen.}
  \label{tab:pilot200}
  \begin{tabular}{lccccc}
    \toprule
    Seed & Model & Params & Val BPB & Val ppl & Train sec \\
    \midrule
    7 & Transformer-SwiGLU & 223,040 & 3.949 & 15.44 & 2.75 \\
    7 & Transformer-GELU & 174,848 & 3.930 & 15.24 & 2.27 \\
    7 & Readout (152) & 221,960 & 3.798 & 13.91 & 32.99 \\
    7 & Stable (184) & 226,004 & 3.832 & 14.24 & 64.84 \\
    11 & Transformer-SwiGLU & 223,040 & 3.946 & 15.41 & 2.23 \\
    11 & Transformer-GELU & 174,848 & 3.965 & 15.62 & 2.00 \\
    11 & Readout (152) & 221,960 & 3.865 & 14.57 & 33.66 \\
    11 & Stable (184) & 226,004 & 3.819 & 14.11 & 64.76 \\
    \bottomrule
  \end{tabular}
\end{table}

These results already show that localized recurrent FFNs can become a
language-modeling win in the small matched-capacity regime. But the gap is
still modest, and the best small-model recurrent designs appear to optimize
more slowly than the Transformer.

\subsection{Accuracy-First Ultra Push}
We next asked a less conservative question: what happens if we optimize purely
for quality inside the FFN slot, even before restoring parameter fairness? The
answer is that widening and compounding the recurrent FFN creates a much larger
jump. Table~\ref{tab:ultra120} shows the most informative 120-step widening
results.

\begin{table}[t]
  \centering
  \caption{Accuracy-first widening results at 120 steps.}
  \label{tab:ultra120}
  \begin{tabular}{lcccc}
    \toprule
    Model & Params & Val BPB & Val ppl & Train sec \\
    \midrule
    Transformer-SwiGLU & 223,040 & 4.404 & 21.17 & 1.59 \\
    Stable (256) & 284,864 & 4.366 & 20.63 & 37.43 \\
    Ultra (128) & 482,432 & 4.366 & 20.62 & 70.85 \\
    Ultra (256) & 888,704 & 4.305 & 19.77 & 68.64 \\
    \bottomrule
  \end{tabular}
\end{table}

The key point is not just that Ultra wins, but that it wins by much more than
the earlier fair recurrent variants. At 120 steps, Ultra-256 nearly matches the
short-budget Transformer-GELU row from the earlier table while clearly beating
the strong small Transformer-SwiGLU baseline.

We then extended Ultra training. In single-seed 200-step runs, Ultra-384
already reaches 3.679 and 3.685 BPB across two seeds. In a targeted 300-step
single-seed search, Ultra-384 falls further to 3.567 BPB, and exploratory
widening to Ultra-448 reaches 3.557 in seed 7. These search results suggest
that the compound recurrent FFN has substantial optimization headroom once the
budget is increased.

\subsection{Matched-Capacity Crossover at 400 Steps}
The most important question is whether the Ultra improvement survives a fairer
capacity-controlled comparison. Table~\ref{tab:ultra400} answers that question.
We compare Ultra-384 against a larger Transformer-SwiGLU with nearly identical
parameter count.

\begin{table}[t]
  \centering
  \caption{Two-seed 400-step matched-capacity comparison.}
  \label{tab:ultra400}
  \begin{tabular}{lccccc}
    \toprule
    Seed & Model & Params & Val BPB & Val ppl & Train sec \\
    \midrule
    7 & Transformer-SwiGLU ($d=160$) & 1,294,880 & 3.474 & 11.11 & 8.51 \\
    7 & Ultra-384 & 1,294,976 & 3.254 & 9.54 & 245.10 \\
    11 & Transformer-SwiGLU ($d=160$) & 1,294,880 & 3.438 & 10.84 & 9.39 \\
    11 & Ultra-384 & 1,294,976 & 3.254 & 9.54 & 222.74 \\
    \bottomrule
  \end{tabular}
\end{table}

This is the strongest result in the paper. Ultra-384 improves BPB by 0.220 in
seed 7 and 0.184 in seed 11 relative to the matched-capacity Transformer. The
two-seed mean is 3.254 for Ultra-384 versus 3.456 for the Transformer. In this
small benchmark, the quality gap is no longer marginal.

We still avoid overselling the result. The Transformer control is strong among
the settings we tested, but our Transformer hyperparameter search is much
smaller than the Ultra search. Even so, the matched-capacity crossover is now
substantial enough that it cannot be dismissed as a tiny short-run fluctuation.

\subsection{Throughput Remains the Main Weakness}
The systems story remains decisively negative in the current implementation.
Table~\ref{tab:throughput} compares the matched-capacity Transformer control
with Ultra-384.

\begin{table}[t]
  \centering
  \caption{Forward-only throughput on Apple MPS for the matched-capacity pair.}
  \label{tab:throughput}
  \begin{tabular}{lcc}
    \toprule
    Sequence length & Transformer-SwiGLU tok/s & Ultra-384 tok/s \\
    \midrule
    64 & 204.1k & 6.6k \\
    128 & 281.1k & 6.6k \\
    256 & 243.8k & 6.7k \\
    \bottomrule
  \end{tabular}
\end{table}

Ultra-384 is therefore a quality-first research result, not a systems
replacement. The throughput is roughly 30--40$\times$ worse in this Python
reference implementation, even though the recurrent curve is flatter across
sequence lengths.

\section{Discussion}
\paragraph{Relation to modern RNNs.}
The paper is motivated by the same broad intuition behind Universal
Transformers, Transformer-to-RNN conversion, RWKV-style models, matrix-state
recurrent models, selective state spaces, and delta-rule sequence models:
dynamic token-conditioned memory can replace some functions usually assigned to
static or attention-heavy components \citep{dehghani2019universal,
kasai2021t2r, orvieto2023lru, rwkv2023, eaglefinch2024, mamba2023,
mamba2ssd2024, retnet2023, yang2024deltanet, griffin2024,
gateddeltanet2024, m2rnn2026}. But our experiment is deliberately narrower.
We do not replace the whole backbone. We ask whether the FFN slot alone can be
made recurrent.

\paragraph{Relation to FFN memory interpretations.}
Our motivation also overlaps with a mechanistic line of work that treats FFN
layers as memory-bearing modules rather than generic token-wise nonlinearities.
Key-value memory interpretations of FFNs \citep{geva2020ffnmem}, knowledge-
neuron analyses \citep{dai2021knowledge}, and direct model-editing results that
target MLP weights \citep{meng2022rome} all support the view that part of a
Transformers' stored knowledge lives in the FFN stack. Our localized
recurrent FFN can be read as shifting some of that storage from purely static
parameters into token-conditioned state, while still keeping the surrounding
attention scaffold intact.

\paragraph{What the current evidence supports.}
The strongest honest claim is now:
\begin{quote}
Localized recurrent FFN replacement is a real small-scale quality win plus a
memory-sensitive win, but the strongest result currently comes from an
accuracy-first compound FFN rather than a strict drop-in systems replacement.
\end{quote}
That is a stronger statement than our earlier ``memory win only'' or ``small
matched win'' interpretations, but still far weaker than a general claim that
recurrent layers replace Transformers everywhere.

\paragraph{Why Ultra likely helps.}
The fair small variants suggest that three ingredients matter: token-conditioned
readout, recurrent dynamics with better stability, and a local branch that
restores FFN-like token mixing. Ultra simply stops forcing those ingredients to
share a tiny budget. It gives the model a readout memory, a stable memory, and
a wider local nonlinear branch, then lets the token mix them dynamically.
Viewed this way, Ultra is less a single clever recurrence rule than a
specialized division of labor inside the FFN slot.

\paragraph{Why the ranking changes with training length.}
The matched small-model screen already hints that good recurrent FFNs optimize
more slowly than the Transformer. The 300- and 400-step Ultra results amplify
that pattern. Very short screens understate what these localized recurrent FFNs
can do once the optimization budget is extended.

\section{Limitations}
This is still a pilot study, not a state-of-the-art language-modeling claim.
\begin{itemize}[leftmargin=1.25em]
  \item The experiments are small: byte-level WikiText-2, short sequences, and
  at most 400 training steps on Apple MPS.
  \item The final matched-capacity result is only a two-seed comparison.
  \item The Ultra result follows a targeted search over recurrent designs and
  schedules. Our larger Transformer control search is smaller and may not be
  fully optimized.
  \item The final comparison matches parameter count, not FLOPs, wall-clock
  budget, or hardware efficiency.
  \item The recurrent implementations are explicit Python loops rather than
  fused scan kernels.
  \item Ultra is not the cleanest possible scientific isolation of a
  parameter-matched FFN replacement. It is explicitly an accuracy-first
  compound design.
  \item We do not yet report larger datasets, matched-FLOP comparisons, or
  modern long-context benchmarks.
  \item The best synthetic memory result is still modest in absolute accuracy,
  even though it improves on the pilot baselines.
\end{itemize}

\section{Conclusion}
Replacing the Transformer FFN with localized recurrent memory is a credible and
scientifically useful research direction. A plain recurrent FFN already helps
memory-sensitive synthetic tasks, which supports the view that the FFN slot can
benefit from a dynamic memory mechanism. A wider recurrent method screen then
shows that several localized recurrent FFNs beat a small strong SwiGLU
Transformer. The strongest update in this paper is that an accuracy-first
compound recurrent FFN, Ultra-384, goes further: in a two-seed 400-step
matched-capacity comparison, it substantially beats a larger Transformer-SwiGLU
on byte-level WikiText-2. The implementation is still drastically slower, so
this should not be framed as a practical replacement yet. But it is now strong
enough to justify a more ambitious standalone line of work: optimize the
compound recurrent FFN itself, add fused kernels, run a stronger Transformer
control search, and test whether the same quality crossover persists on larger
datasets.

\bibliographystyle{plainnat}
\bibliography{refs}

\end{document}
